% ===========================================================================
%  user/09-account/00-chapter.tex
% ===========================================================================
\chapter{Account and subscription}\label{cap:account}

\section{What the account is for}\label{sez:account-for}

Four things, and nothing else: entering the marketplace and receiving
messages --- both in the \voce{Community} --- having a public profile, and
taking out a subscription. The rest of the application works without one.

\figuraSchermata{0.86}{figure/schermate/subscription.png}
  {Pages that need an account say so and offer to sign you in, instead of
   disappearing from the menu.}
  {signed-out}

\section{Signing in for the first time}\label{sez:first-sign-in}

\begin{pratica}{During the beta}
During the beta you do not create an account yourself: Lyra creates it when
you ask to join, on the site, with \voce{Entra nella beta}.

\begin{enumerate}
  \item When the beta starts, an email arrives with the link to download the
    application, your address and a \textbf{first password}.
  \item Install the application, press \pulsante{Sign in} and enter that
    address and that password.
  \item The application asks you to choose your own password, twice. From
    then on you sign in with it: the one in the email no longer works.
\end{enumerate}

If you already had an account with that address, you keep it and its
password: the email only brings the link.
\end{pratica}

\figuraSchermata{0.86}{figure/schermate/beta-first-password.png}
  {The first sign-in with the password from the invitation: the application
   asks for your own, at least eight characters and different from the one
   in the email.}
  {beta-first-password}

\noindent
The places in the beta are limited. When they are full, a request stays on
the waiting list: the site says so, and if a place frees up the team adds you
and the email arrives. An address that did not ask to join is refused, and so is an
attempt to sign up from the application while the beta lasts; the sign-in
page says which of these is the case.

\noindent
You can also sign in with Google (\voce{Continue with Google}), with the
address of an account you already have: the sign-in opens in the system
browser and the application picks up the result.

\section{The profile}\label{sez:profile}

\voce{Profile}, in the account menu at the top right, gathers the account on
one page. \pulsante{Sign out} is at the top right of the page.

\figuraSchermata{1.0}{figure/schermate/profile.png}
  {The top of the profile: the plan, two-step verification, the payout account
   and the published artifacts at a glance; below, the account details,
   marketplace publishing and earnings.}
  {profile}

\begin{description}[leftmargin = 1.4em, style = nextline, itemsep = 4pt]
  \item[\voce{Account details}] Display name and bio --- what others see next
    to a published artifact --- changed with \pulsantesec{Edit profile}.
  \item[\voce{Marketplace publishing}] Whether you can sell, with links to
    your purchases and your listings, and to \voce{My subscription}.
  \item[\voce{Marketplace earnings}] The payout account that receives what
    you sell. Selling for a price needs one, and linking it needs two-step
    verification first.
  \item[\voce{Automatic publishing}] Off unless you turn it on
    (\vediSez{auto-publish}).
  \item[\voce{Security}] The password, the email address and two-step
    verification. They are not stored by the application: they live on the
    site, and the application only passes them on.
  \item[\voce{Preferences}, \voce{Brokers and data keys}] The theme, how
    pages open, hardware acceleration and the notifications --- the same
    choices as the \voce{Settings} window (\vediSez{speed}); the broker
    connections and the data keys, which stay on this computer.
  \item[\voce{Send us feedback}] A report or a suggestion for the Lyra
    team.
\end{description}

\subsection{Automatic publishing}\label{sez:auto-publish}

Lyra can upload your best work to the marketplace by itself --- only if you
ask it to. With the switch on, whenever one of your strategies, approaches
or projects reaches an out-of-sample walk-forward Sharpe above 2.0, Lyra
uploads it under your account, code and results, \textbf{free and not
public}: only the Lyra administrators see it. Nothing gets a price, nothing
appears in the storefront, nothing is shown to other users.

\figuraSchermata{0.56}{figure/schermate/profile-autopublish.png}
  {Automatic publishing in the profile, switched on. Until you turn it on
   the line reads \enquote{Off: nothing is uploaded.}}
  {profile-autopublish}

\begin{itemize}
  \item The choice is for the account you are signed in with, on this
    computer. It is off for everyone who has not turned it on --- including
    after an update.
  \item You can turn it off at any time, from the same switch. What was
    already uploaded is listed under the switch, and can be removed from the
    marketplace.
  \item What you bought in the marketplace is never uploaded: it is someone
    else's work.
\end{itemize}

\begin{nota}{How Lyra recognises this installation}
With a random number, made the first time Lyra starts and kept with your
data. It says nothing about your computer: before, it was derived from the
computer's hardware, network card included. It counts towards the limit of
installations of a purchase; deleting it makes Lyra a new installation.
\end{nota}

\section{Speed: how pages open, hardware acceleration}\label{sez:speed}

Two choices in the \voce{Settings} window (account menu at the top right),
and in the preferences of the profile.

\begin{description}[leftmargin = 1.4em, style = nextline, itemsep = 4pt]
  \item[\voce{Opening pages}] \voce{Right away}, the default: a new page
    opens at once and fills in as its data arrives, with the thin bar at the
    top of the window running until it has everything. \voce{When loaded}:
    the page you are leaving stays on screen until the new one has its data
    and charts, up to two seconds --- no page filling in before your eyes,
    steadier on slower computers. Quick pages look the same either way. The
    choice applies from the next page you open.
  \item[\voce{Hardware acceleration}] Only in the desktop application; in a
    browser it is a setting of the browser. \voce{Automatic}, the default,
    lets the graphics card draw the pages whenever Chromium supports it:
    almost always the fastest. \voce{Always} uses the card even when Chromium
    does not trust it (older integrated graphics, old drivers): smoother
    pages where the processor was drawing them, but a faulty driver can make
    the window flicker or go black. \voce{Off} draws with the processor only,
    for driver problems; it is slower. Below the buttons, what the card is
    doing in this session. A new choice applies when the application
    restarts: \pulsantesec{Restart now} does it.
\end{description}

\begin{attenzione}{If the window goes black after choosing \voce{Always}}
Start the application with the environment variable \cod{LYRA\_GPU=off}
(or \cod{auto}): it overrides the saved choice for that session, and from
\voce{Settings} you can set \voce{Automatic} again. If the graphics card
stops responding twice with \voce{Always}, the application goes back to
\voce{Automatic} by itself at the next start, and the section says so.
\end{attenzione}

\section{The plans}\label{sez:plans}

The plans are the same as on the Pricing page of the site. Free is the whole
application; you pay for privacy, data and a higher share of what your
strategies earn --- not for using the application. They are on the
\voce{My subscription} page, in the account menu at the top right.

\begin{center}
\begin{tabular}{@{}L{22mm}L{30mm}L{86mm}@{}}
\toprule
\textbf{Plan} & \textbf{Price} & \textbf{What it adds} \\
\midrule
\rigaTenue
\voce{Free} & \eur{0}, forever & The whole application: building and
  backtesting locally with your data, the local AI copilot, out-of-sample
  validation, live runners and kill-switch. On the marketplace: browse,
  download, buy. Share of earnings \perc{0}: an allocated strategy enters the
  track record. \\
\voce{Pro} & \eur{49} a month or \eur{490} a year & Full privacy for code
  and metrics; selling on the marketplace and buying at $-$\perc{10}; the
  institutional-data sandbox (from 2030); \perc{10} of what your allocated
  strategies earn. \\
\rigaTenue
\voce{Professional} & \eur{149} a month or \eur{1,490} a year & Raw history
  on your machine (from 2030) with $-$\perc{10} on data vendors; marketplace
  $-$\perc{20}; sandbox on more markets; \perc{30} of earnings. It is the
  recommended plan. \\
\voce{Quant} & from \eur{200} a month or \eur{2,000} a year, plus data &
  Exchange-licensed real-time (from 2032), in per-market packages; extended
  sandbox; \perc{30} of earnings. \\
\rigaTenue
\voce{Desk} & \eur{8,000} a year & Five seats, installation on your servers,
  compliance export, white-label and SLA; \perc{30} of earnings, per seat, by
  contract. Sold by contract, after a demo. \\
\bottomrule
\end{tabular}
\end{center}

\figuraSchermata{1.0}{figure/schermate/subscription-plans.png}
  {The plans on \voce{My subscription}: price, share of earnings and what
   each plan adds. \voce{Show the detailed comparison} opens the table with
   every feature, plan by plan.}
  {subscription-plans}

\begin{nota}{The features with a year}
Some features arrive with a later release: the sandbox and the raw history
in 2030, real-time in 2032. The cards and the comparison table show the year
next to them.
\end{nota}

\begin{nota}{During the beta}
The paid plans open after the beta. Until then Free is the whole
application, nothing can be bought or sold on the marketplace, nothing is
charged, and the plan buttons say \voce{On sale after the beta}. Desk demo
requests open later too: until then the Desk button says \voce{Coming soon}.
\end{nota}

\begin{nota}{The share of earnings}
It is the part of what a strategy earns, once it is allocated to capital,
that goes to its author: \perc{10} with Pro, \perc{30} with Professional,
Quant and Desk, set by contract. On Free an allocated strategy enters the
track record without a share. Every submission is opt-in, strategy by
strategy and under a contract: nothing is submitted without your consent.
The site describes the steps, from the protocol to the live track record.
\end{nota}

\subsection{The two-month trial}

Pro, Professional and Quant start with a two-month trial, one per account.

\begin{pratica}{Starting the trial}
\begin{enumerate}
  \item On the plan's card, press \pulsante{Start 2-month trial}; the hint
        below the button says until when nothing is charged.
  \item The window shows \voce{Charged today} at \eur{0}, the day the trial
        ends and the price that applies afterwards. Choose monthly or yearly
        there, then press \pulsante{Start the trial}.
  \item The Stripe page opens: Stripe asks for your card but charges
        nothing. The page updates by itself and shows \voce{Trial until} and
        the date.
  \item On that date the first payment is taken. Cancel before, and nothing
        is ever charged.
\end{enumerate}
\end{pratica}

\noindent
Changing plan during the trial costs nothing, and the trial keeps its end
date. An account that has already had a trial, or a paid plan, pays from the
first day: the button then says \pulsante{Choose} followed by the plan's
name.

\subsection{Desk}

Desk is not bought with a card: it starts with a demo on your firm's own
strategies, and a contract. Demo requests for firms are not open yet. Until
they are, \pulsantesec{Request a demo} on the Desk card is switched off:
point at it, or reach it with \voce{Tab}, and it says
\voce{Coming soon}; under it, \voce{Sold by contract · demos open soon}. The
Desk buttons on the site's Pricing page say the same. When requests open, the
button opens a short form --- the firm, how many people would use the
platform, a message --- and the request reaches the Lyra team, who reply by
email and in \voce{Messages}.

\subsection{Monthly or yearly}

Both. Yearly costs ten monthly payments: two months are free. Desk is yearly
only. You choose when you subscribe and switch whenever you like, from the
button next to your plan: Stripe works out the adjustment for the period
already paid.

\subsection{Changing plan}

\pulsantesec{Upgrade} and \pulsantesec{Downgrade} on the cards change the plan
at once: Stripe credits what you have already paid for the current period and
charges the difference to the card on file. Before you confirm, the window
lists what you gain or what you lose.

\subsection{Cancelling, card, invoices}

From \pulsante{Manage in Stripe} on \voce{My subscription}. The billing
portal hosted by Stripe opens: there you cancel the subscription, change the
payment method and download the invoices. The card never reaches the
application.

\begin{nota}{Cancelling takes nothing you already have}
The plan stays active until the end of the period you paid for. A purchased
artifact stays purchased and installed, and strategies, code and data stay on
your computer: you go back to Free, which is still the complete application.
What stops is selling on the marketplace and the benefits of the plan.
\end{nota}

\section{When something does not add up}\label{sez:account-problems}

\begin{center}
\begin{tabular}{@{}L{50mm}L{88mm}@{}}
\toprule
\textbf{Symptom} & \textbf{What to do} \\
\midrule
\rigaTenue
\enquote{You cannot sell} but you have a subscription
  & Close and reopen the application: the link with the site is made at
    start-up. \\
The marketplace is empty
  & Same reason. If it persists, sign out and sign back in. \\
\rigaTenue
I paid and I cannot see the artifact
  & Wait a few seconds and reload. If it stays like that, use the command
    that explains why it does not show as purchased, on the artifact's page
    (\vediCap{community}). \\
The trial is not offered
  & The trial is one per account: an account that has already had one, or a
    paid plan, pays from the first day. \\
\rigaTenue
A page says it does not know, instead of showing a number
  & Not a fault: the application could not ask the site, and says so instead
    of making up an answer. Try again later. \\
\bottomrule
\end{tabular}
\end{center}
